\documentclass[12pt,fleqn]{report}

\voffset=-06mm
\hoffset=-16mm
\textwidth=160mm
\textheight=225mm

\setcounter{totalnumber}{5}
\renewcommand{\textfraction}{0.1}
\parskip=8pt plus2pt minus2pt

\include{MHEADS}

\begin{document}

\tableofcontents

\newpage

\pagenumbering{arabic}

\chapter{Queueing theory}

\begin{itemize}

\item Queueing theory enables us to check the viability of
complicated systems, i.e. given a hardware configuration 
and a workload, we can predict job throughput, hardware 
utilisations and bottlenecks.

\item Applications include operating systems,
computer networks, industrial complexes, ...

\item There are a number of standard patterns with
known solutions. Your job is to classify the problem 
into one of these patterns, where you can.
If you cannot, then you may have to resort to simulation.

\end{itemize}

In these notes, we look at some of the basic
standard patterns, and apply them to some simple
computer systems and networks.


\noindent
Queueing theorists traditionally use three Greek letters:


\begin{center}
\begin{tabular}{l c l c l}
$\lambda$ && lambda && arrival rate \\[0.2in]
$\mu$     && mu     && service rate \\[0.2in]
$\rho$    && rho    && utilisation (how busy) \\[0.2in]
\end{tabular}
\end{center}


\section{Preliminary}
First two reminders of how to work out some well-known 
and elementary formulae used in queueing theory.

\subsection{1/1-x}

To find the sum of 1 + x + x$^2$ + x$^3$ + ....
Its not obvious so let's call it S. Then, by definition,

\begin{tabular}{r cc cc cc cc ccc}
S & = & 1 & + & x & + & x$^2$ & + & x$^3$ & + & x$^4$ & .... \\
\end{tabular}

\noindent
We notice that each term on the right is greater by x than the one
before.
Hmmh. This suggests we look at  S.x (S multiplied by x), 
more conveniently written the other way round as x.S.

\begin{tabular}{r cc cc cc cc ccc}
x.S & = & x & + & x$^2$ & + & x$^3$ & + & x$^4$ &  .... \\
\end{tabular}

\noindent
The right hand sides of S and x.S are almost the same
(its helpful to ``get the columns right'').


\begin{tabular}{r cc cc cc cc ccc}
S   & = & 1 & + & x & + & x$^2$ & + & x$^3$ & + & x$^4$ & .... \\
x.S & = &   &   & x & + & x$^2$ & + & x$^3$ & + & x$^4$ &  .... \\
\end{tabular}


\noindent
What if we subtract? 

\begin{tabular}{r cc cc cc cc ccc c}
S   & = & 1 & + & x & + & x$^2$ & + & x$^3$ & + & x$^4$ & .... \\
x.S & = &   &   & x & + & x$^2$ & + & x$^3$ & + & x$^4$ &  .... \\
\hline
S-x.S & = & 1 & \\
\hline
\end{tabular}

\noindent
Hence  S-x.S = 1, or S = 1/(1-x). So that we have shown that

\begin{table}[hbt]\centering
\begin{tabular}{||r cc cc cc cc ccc||}
\hline
&&&&&&&&&&&\\
\hspace*{0.25in}
1/(1-x)
  & = & 
1 & + & x & + & x$^2$ & + & x$^3$ & + & x$^4$ & .... 
\hspace*{0.25in}\\
&&&&&&&&&&&\\
\hline
\end{tabular}
\end{table}

\subsection{1/(1-x)$^2$}

We also need to recognise the expansion of $\frac{1}{(1-x)^2}$.
We start off with the obvious fact that $\frac{1}{(1-x)^2}$
= $\frac{1}{(1-x)}$ $\times$ $\frac{1}{(1-x)}$, or S $\times$ S,
or as we have just shown 

\begin{center}
(1 + x + x$^2$ + x$^3$ + ....) $\times$ (1 + x + x$^2$ + x$^3$ + ....)
\end{center}

\noindent
the best way to proceed is to break the problem up into smaller steps


\begin{tabular}{r cc cc cc cc cc cc cc}
\multicolumn{10}{l}{$\frac{1}{(1-x)^2}$} \\
=  & 1 & $\times$ &(& 1 & + & x & + & x$^2$ & + & x$^3$ & + & x$^4$ &
  .... &)\\
+  & x & $\times$ &(& 1 & + & x & + & x$^2$ & + & x$^3$ & + & x$^4$ &
  .... &)\\
+  & x$^2$ & $\times$ &(& 1 & + & x & + & x$^2$ & + & x$^3$ & + &
  x$^4$ & .... &)\\
+  & x$^3$ & $\times$ &(& 1 & + & x & + & x$^2$ & + & x$^3$ & + & x$^4$ & .... &)\\
+  & x$^4$ & $\times$ &(& 1 & + & x & + & x$^2$ & + & x$^3$ & + &
  x$^4$ & .... &)\\
+ & \multicolumn{10}{l}{.........................................}\\
\end{tabular}

\noindent
Carrying through the multiplications line by line and taking care to
align so that we have columns of the same power, we get 

\begin{tabular}{r cc cc cc cc cc cc cc}
\multicolumn{10}{l}{$\frac{1}{(1-x)^2}$} \\
=  & 1 & + & x & + & x$^2$ & + & x$^3$ & + & x$^4$ &  .... &\\
   &   & + & x & + & x$^2$ & + & x$^3$ & + & x$^4$ &  .... &\\
   &   &   &   & + & x$^2$ & + & x$^3$ & + & x$^4$ &  .... &\\
   &   &   &   &   &       & + & x$^3$ & + & x$^4$ &  .... &\\
   &   &   &   &   &       &   &       & + & x$^4$ &  .... &\\
\multicolumn{10}{l}{.........................................}\\
\hline
= &  1 & + & 2.x & + & 3.x$^2$ & + & 4.x$^3$ & + & 5.x$^4$ &  .... &\\
\hline
\end{tabular}

\noindent
So that we have shown that

\begin{table}[hbt]\centering
\begin{tabular}{||r cc cc cc cc ccc||}
\hline
&&&&&&&&&&&\\
\hspace*{0.25in}
1/(1-x)$^2$
  & = & 
       1 & + & 2.x & + & 3.x$^2$ & + & 4.x$^3$ & + & 5.x$^4$ & .... 
\hspace*{0.25in}\\
&&&&&&&&&&&\\
\hline
\end{tabular}
\end{table}

\newpage

\section{Introduction}

\subsection{Intuition}

For intuition, consider a small shop with one assistant
and three customers who arrive at times 1, 2, and 10.
The customers are served on the first-come, first-served basis
and each service takes 3 time units (say minutes).
The assistant can only serve one customer at a time.

Here is a trace of the significant times as our system unfolds
(for 15 minutes).

{\small 
\begin{center}
\begin{tabular}{|| r | l | l | c|| }
\hline
Time & Event      &                  & Server \\
\hline
0    &            &                  & Idle   \\
1    & C1 arrives & service 1 starts & Busy \\ 
2    & C2 arrives & must wait        & Busy \\
3    &            &                  & Busy \\
4    &            & service 1 ends   & Busy  \\
     &            & C1 exits         & \\
     &            & service 2 starts & \\
5    &            &                  & Busy \\
6    &            &                  & Busy \\
7    &            & service 2 ends   & Idle  \\
     &            & C2 exits         & \\
8    &            &                  & Idle \\
9    &            &                  & Idle \\
10   & C3 arrives & service 1 starts & Busy \\ 
11   &            &                  & Busy \\
12    &            &                  & Busy \\
13    &            & service 3 ends   & Idle  \\
      &            & C3 exits         & \\
14    &            &                  & Idle \\
15    &            & shutdown         &      \\
\hline
\end{tabular}
\end{center}}



In 15 minutes, the server is busy for 9 minutes and idle for 
6 minutes.
The server utilisation, $\rho$, = 9/15, i.e. 0.66.. or 66\%.

The service rate, $\mu$, the number of customers servable per minute,
is the reciprocal of the expected (average) service time of the
server. $\mu$ = 1/3 = 0.33 customers per minute.

In 15 minutes there are three arrivals. The arrival rate $\lambda$ =
3/15 = 0.2 customers per minute.
In 15 minutes there are 3 departures. 
The departure rate = 3/15 = 0.2 customers per minute.
Note that in a viable system (ones in which queues do
not build forever, what comes in, goes out, so the
departure rate will equal the arrival rate.

Notice that $\rho$ = 0.66 = $\frac{\lambda}{\mu}$.

Notice further that the expected through time (time a customner spends
in the system = (3 + 4 + 3)/3 = 3.33 minutes.
And that the expected time spent queueing is (0 + 1 + 0)/3 = 0.33
minutes.
Here the expected through time = 3.33 = expected service time = expected
wait. Is this a reasonable thing to expect?

\subsection{More serious}

Usually some properties of a queueing system are known,
and we have to deduce appropriate measures of system 
performance from them.
Typically we will be given 

\begin{tabular}{r c l}
$\lambda$ &  & the expected job (customer, message, ...) arrival rate \\
$\mu$     &  & the expected service rate per server           \\
\end{tabular}

\noindent
If on average, 20 jobs arrive every 4 minutes, 
then $\lambda$, the arrival rate is 5 jobs/min (or 300 jobs/hour).
It is sometimes useful to refer to the expected (or average)
inter-arrival time, denoted by E[$\tau$], which in this case is 
12 seconds (or 0.2 minutes, or $\frac{1}{300}$ hours).

Similarly, if there is one server whose service rate is 
20 jobs every 2 minutes, then $\mu$ is 10 jobs/min (or 600 jobs/hour).
In this case, the expected (or average) service time, denoted by E[s],
is 6 seconds (or 0.1 minutes, or $\frac{1}{600}$ hours).

One measure of performance is {\it server utilisation}
 $\rho$ = $\lambda$/$\mu$ =  $\lambda$.E[s]

Using the data above, $\rho$ =  0.5.
$\lambda$ and $\mu$ are essentially positive.
It spells trouble if jobs arrive faster than we can serve them,
so to make a system viable, we will normally want $\rho$ $<$ 1.
If it isn't then we will have to make the server faster,
supply more servers, or limit the queue size and
discard jobs we cannot handle.

\subsection*{Example 1}

 A stream of short jobs arrives at a cpu for execution
to completion. What is the average utilisation if
given :

(i) an arrival rate of 10 jobs per minute 
         and a mean service time of 4 seconds per job, and

(ii) an arrival rate of 12 jobs per minute 
         and a mean service time of 4 seconds per job.

(iii) an arrival rate of 20 jobs per minute 
         and a mean service time of 4 seconds per job.

\noindent
{\bf Answer}.

(i) \begin{tabular}{r c l r}
     $\lambda$ & = & 10 jobs/min \\
     $\mu$     & = & 15 jobs/min \\
     $\rho$    & = & 0.666       & $\rho$ $<$ 1, i.e. the system is
     viable \\
     \end{tabular}

(ii)      \begin{tabular}{r c l r}
     $\lambda$ & = & 12 jobs/min \\
     $\mu$     & = & 15 jobs/min \\
     $\rho$    & = & 0.800       & $\rho$ $<$ 1, i.e. the system is
     viable \\
     \end{tabular}

(iii)      \begin{tabular}{r c l r}
     $\lambda$ & = & 20 jobs/min \\
     $\mu$     & = & 15 jobs/min \\
               &    &           &  $\lambda$ $>$ $\mu$ and the system is not
     viable \\
     \end{tabular}

\noindent
\rule{160mm}{0.5mm} 

Besides $\rho$, other useful measures of system performance are

\begin{tabular}{r c l}
T     & & the average waiting time in the system \\
T$_q$ & & the average waiting time in the queue  \\
N     & & the average number of jobs in the system \\
N$_q$ & & the average number of jobs in the queue \\
$\pi$(90) &  & the 90th percentile of waiting times in the system\\
$\pi$$_w$(90) &  & the 90th percentile of waiting times in the queue\\
P$_n$ & & the probability of n jobs in the system  \\
\end{tabular}


Little has shown that for a steady state queueing system, under very
general conditions, then:

\begin{center}
\begin{tabular} { || r c l ||}
\hline
\multicolumn{1}{||l}{\bf Little's Law}     &&\\
   &   &            \\
\hspace*{1in} N  & = & $\lambda$T \hspace*{1in} \\
   &   &       \\
N$_q$  & = & $\lambda$T$_q$ \\
   &   &            \\
    \hline
\end{tabular}
\end{center}

\noindent
Given  $\lambda$ and E[s] (or 1/$\mu$), Little's Law
and the fact that T = T$_q$ + E[s] (the average 
through time is the sum of the average time spent queueing plus 
the expected  service time) enables us
to compute four of our key performance measures if we know one of
them. For example if T is known, then

    
\begin{tabular}{ r c l r }
N      & = &  $\lambda$T & Little's law \\
&&&\\
T$_q$  & = & T -- E[s]  & since T = T$_q$ + E[s] \\
&&&\\
N$_q$  & = & $\lambda$T$_q$ & Little's Law again\\
\end{tabular}

\noindent
The computation of the four measures varies with the queueing model.

\subsection{The exponential nature of random arrivals}

Throughout our look at systems of queues we are assuming that
our arrival and service patterns do not depend upon
the state of the system (i.e. they remain the same whether the system
is heavily loaded or lightly loaded).

We start off with the reasonable assumption that the number of
arrivals in a time slot (period of observation) 
depends upon the length of the time slot.
 
\begin{tabular}{l c l}
\rule{1in}{0.05in} & & expect n arrivals \\
\rule{2in}{0.05in} & & expect 2n arrivals \\
\rule{3in}{0.05in} & & expect 3n arrivals \\
\multicolumn{3}{c}{.............}\\
\end{tabular}

\noindent
If the arrival rate is $\lambda$ then we expect 
to average $\lambda$.t arrivals in a time slot of width t.

Now let's look at the probability of NO arrivals in a time slot 
of width t. We call this probability P(t), and
by definition the probability of no arrivals in a time slot 
of width t+$\delta$t is P(t+$\delta$t).

Now we get no arrivals in a time slot of width 
t+$\delta$t by having no arrival in the slot of width t
(and the probability of this happening is P(t))
and having no arrival in the next period of length $\delta$t.
And since the probability of an arrival in a period of length $\delta$t
is $\lambda$.$\delta$t (our starting assumption), 
then the probability of no arrival
in the period $\delta$t is (1- $\lambda$.$\delta$t).
So that 

\begin{center}
P(t+$\delta$t) = P(t) $\times$ (1 - $\lambda$.$\delta$t)
\end{center}

But the Taylor expansion of P(t+$\delta$t) is

\begin{center}
P(t+$\delta$t)  =  P(t)  +  $\delta$t.P'(t) + {\Large O}($\delta$t$^2$) ....
\end{center}

Combining these last two formulae we get

\begin{center}
P(t)  +  $\delta$t.P'(t) .... =  P(t) -- $\lambda$.$\delta$t.P(t)
\end{center}


As $\delta$t gets vanishingly small, we may neglect second order terms,
so 

\begin{center}
$\delta$t.P'(t) = -- $\lambda$.$\delta$t.P(t)
\end{center}

i.e. P'(t) = -- $\lambda$.P(t), the solution to which is 


\begin{table}[hbt]\centering
\begin{tabular}{|| cc c||}
\hline
&&\\
\hspace*{0.25in} P(t)   & = &  e$^{-\lambda t}$ \hspace*{0.25in}\\
&&\\
\hline
\end{tabular}
\end{table}

Remember that P(t) is is the probability of no arrivals in 
a time period of length t, so the probability of 
some arrivals in this time period is (1 -- e$^{-\lambda t}$).

\newpage

\section{M/M/1 - the single server queue}

The simplest queueing model is the single server queue
with customers taken FCFS. If the system is not empty,
then the job being served is NOT considered to be part of the queue.
The queue is as long as need be (potentially infinite).

We start our analysis by characterizing the possible states the system
can be in using the number of jobs in the system as our measure.
There can be 0 jobs in the system, 1 job, 2 jobs, ....
We designate these states  S$_0$, S$_1$, S$_2$, ...

\subsection{State transitions}

\begin{figure}[hbt]\centering
\setlength{\unitlength}{0.20mm}
\begin{picture}(720,240)(0, 0)

\thicklines

\put(   0,   0){\framebox(720, 240){}}

\put(120,  120){\oval(60, 60)} 
\put(120,  120){\makebox(0, 0){\sf 0}}


\put(240,  120){\oval(60, 60)} 
\put(240,  120){\makebox(0, 0){\sf 1}}


\put(360,  120){\oval(60, 60)} 
\put(360,  120){\makebox(0, 0){\sf 2}}


\put(480,  120){\oval(60, 60)} 
\put(480,  120){\makebox(0, 0){\sf 3}}


\put(600,  120){\makebox(0, 0){\sf ...................}}


\end{picture}
\label{M/M/1-states}
\caption{System States}
\end{figure}




How can we make a transition  from one state to another?

Consider state S$_0$, the state in which there are no jobs in the system.
The server is idle and the queue is empty.
All that can happen is a job arrives 
and then we flip from state S$_0$ to state S$_1$.
This is modelled below with the transition between states marked
by the rate at which it can occur.

\begin{figure}[hbt]\centering
\setlength{\unitlength}{0.20mm}
\begin{picture}(720,240)(0, 0)

\thicklines

\put(   0,   0){\framebox(720, 240){}}

\put(120,  120){\oval(60, 60)} 
\put(120,  120){\makebox(0, 0){\sf 0}}
\put(150,  140){\vector(1, 0){60}} 
\put(180,  150){\makebox(0, 0){$\lambda$}}


\put(240,  120){\oval(60, 60)} 
\put(240,  120){\makebox(0, 0){\sf 1}}


\put(360,  120){\oval(60, 60)} 
\put(360,  120){\makebox(0, 0){\sf 2}}


\put(480,  120){\oval(60, 60)} 
\put(480,  120){\makebox(0, 0){\sf 3}}


\put(600,  120){\makebox(0, 0){\sf ...................}}


\end{picture}
\label{M/M/1-0}
\caption{S0}
\end{figure}



Now consider state S$_1$, the state in which the server is busy
and the queue is empty.
All that can happen is a fresh arrival (that job will start the queue)
which moves us to state S$_2$,
or the service completes which moves us back to an empty system (state
S$_0$).

\begin{figure}[hbt]\centering
\setlength{\unitlength}{0.20mm}
\begin{picture}(720,240)(0, 0)

\thicklines

\put(   0,   0){\framebox(720, 240){}}

\put(120,  120){\oval(60, 60)} 
\put(120,  120){\makebox(0, 0){\sf 0}}

\put(240,  120){\oval(60, 60)} 
\put(240,  120){\makebox(0, 0){\sf 1}}
\put(270,  140){\vector(1, 0){60}} 
\put(300,  150){\makebox(0, 0){$\lambda$}}
\put(210,  100){\vector(-1, 0){60}} 
\put(180,   90){\makebox(0, 0){$\mu$}}

\put(360,  120){\oval(60, 60)} 
\put(360,  120){\makebox(0, 0){\sf 2}}


\put(480,  120){\oval(60, 60)} 
\put(480,  120){\makebox(0, 0){\sf 3}}


\put(600,  120){\makebox(0, 0){\sf ...................}}


\end{picture}
\label{M/M/1-1}
\caption{S1}
\end{figure}


Now consider state S$_2$, the state in which the server is busy
and the queue has length 1.
All that can happen is a fresh arrival (that job will join the queue)
which moves us to state S$_3$,
or the service completes at which time the current job leaves the
system and the (next) job in the queue starts being served
which takes us back to state S$_1$.


\begin{figure}[hbt]\centering
\setlength{\unitlength}{0.20mm}
\begin{picture}(720,240)(0, 0)

\thicklines

\put(   0,   0){\framebox(720, 240){}}

\put(120,  120){\oval(60, 60)} 
\put(120,  120){\makebox(0, 0){\sf 0}}

\put(240,  120){\oval(60, 60)} 
\put(240,  120){\makebox(0, 0){\sf 1}}

\put(360,  120){\oval(60, 60)} 
\put(360,  120){\makebox(0, 0){\sf 2}}
\put(390,  140){\vector(1, 0){60}} 
\put(420,  150){\makebox(0, 0){$\lambda$}}
\put(330,  100){\vector(-1, 0){60}} 
\put(300,   90){\makebox(0, 0){$\mu$}}



\put(480,  120){\oval(60, 60)} 
\put(480,  120){\makebox(0, 0){\sf 3}}


\put(600,  120){\makebox(0, 0){\sf ...................}}


\end{picture}
\label{M/M/1-2}
\caption{S2}
\end{figure}


Etc, etc. We can now piece together the state transition diagram for
the whole system (or as much as we wish to draw!).

\begin{figure}[hbt]\centering
\setlength{\unitlength}{0.20mm}
\begin{picture}(720,240)(0, 0)

\thicklines

\put(   0,   0){\framebox(720, 240){}}

\put(120,  120){\oval(60, 60)} 
\put(120,  120){\makebox(0, 0){\sf 0}}
\put(150,  140){\vector(1, 0){60}} 
\put(180,  150){\makebox(0, 0){$\lambda$}}


\put(240,  120){\oval(60, 60)} 
\put(240,  120){\makebox(0, 0){\sf 1}}
\put(270,  140){\vector(1, 0){60}} 
\put(300,  150){\makebox(0, 0){$\lambda$}}
\put(210,  100){\vector(-1, 0){60}} 
\put(180,   90){\makebox(0, 0){$\mu$}}

\put(360,  120){\oval(60, 60)} 
\put(360,  120){\makebox(0, 0){\sf 2}}
\put(390,  140){\vector(1, 0){60}} 
\put(420,  150){\makebox(0, 0){$\lambda$}}
\put(330,  100){\vector(-1, 0){60}} 
\put(300,   90){\makebox(0, 0){$\mu$}}

\put(480,  120){\oval(60, 60)} 
\put(480,  120){\makebox(0, 0){\sf 3}}
\put(510,  140){\vector(1, 0){60}} 
\put(540,  150){\makebox(0, 0){$\lambda$}}
\put(450,  100){\vector(-1, 0){60}} 
\put(420,   90){\makebox(0, 0){$\mu$}}

\put(600,  120){\makebox(0, 0){\sf ...................}}

%\put(630,  140){\vector(1, 0){60}} 
%\put(660,  150){\makebox(0, 0){$\lambda$}}
\put(570,  100){\vector(-1, 0){60}} 
\put(540,   90){\makebox(0, 0){$\mu$}}

\end{picture}
\label{M/M/1-n}
\caption{M/M/1 transitions}
\end{figure}

\subsection{Equations of flow balance}


The idea now is to look at equations of balance between states.
Initially the system is empty and it will take a while to ``warm up''. 
Once things have settled down and we have reached 
a ``steady state'', then the average flow between states must balance
(what comes in must go out).
What are these flows? 

Consider the flow between S$_0$ and S$_1$.
There is a $\lambda$ transition to the right and
a $\mu$ transition to the left.
But we must also take into account the probability of being in 
a particular state too.
We use the notation P$_k$ to denote the probability of being in state
S$_k$.
Clearly

\begin{center}
 1 = $\sum_{k=0}^{\infty}$P$_k$ 
    =  P$_0$ + P$_1$ + P$_2$ + P$_3$ + ....
\end{center}

\noindent
Then 

\begin{tabular}{r c l r}
$\lambda$.P$_0$ & = & $\mu$.P$_1$ & flow balance between S$_0$ and
S$_1$ \\
&&&\\
$\lambda$.P$_1$ & = & $\mu$.P$_2$ & flow balance between S$_1$ and
S$_2$ \\
&&&\\
$\lambda$.P$_2$ & = & $\mu$.P$_3$ & flow balance between S$_2$ and
S$_3$ \\
\multicolumn{3}{c}{.................}\\
&&&\\
\end{tabular}

\noindent
Or more neatly

\begin{center}
P$_1$ = $\rho$.P$_0$ \hspace*{0.25in}
P$_2$ = $\rho$.P$_1$ = $\rho^2$.P$_0$ \hspace*{0.25in}
P$_3$ = $\rho$.P$_2$ = $\rho^3$.P$_0$ \hspace*{0.25in}
.....................
\end{center}


\noindent
Thus we have P$_k$ = $\rho^k$P$_0$ (for all k $\geq$ 0):

\begin{tabular}{ r c l}
1 & = & P$_0$ + P$_1$ + P$_2$ + P$_3$ + ....  \\
  & = & P$_0$ + $\rho$.P$_0$ + $\rho^2$.P$_0$ + $\rho^3$.P$_0$ + ....  \\
  & = & P$_0$.(1 + $\rho$ + $\rho^2$ + $\rho^3$ + ....)\\
  & = & P$_0$/(1 - $\rho$)  \hfill for $\rho$ $<$ 1  \\
  & & \\
\end{tabular}

\noindent
>From which 

\begin{tabular}{r c l}
P$_0$ & = & 1 - $\rho$ \hspace*{2in} (probability of an empty system)\\
P$_1$ & = & $\rho$(1 - $\rho$) \\
P$_2$ & = & $\rho^2$(1 - $\rho$) \\
      & {............} & \\
P$_k$ & = & $\rho^k$(1 - $\rho$) \\
\end{tabular}

\subsection{Derivation of key system properties}

N --- the average number of jobs in the system --- 
is computed by:

\begin{tabular}{c c l}
N & = & 0.P$_0$ + 1.P$_1$ + 2.P$_2$ + 3.P$_3$ + ... \\
  & = & 1.$\rho$.P$_0$ + 2.$\rho^2$.P$_0$ + 3.$\rho^3$.P$_0$ + ...\\
  & = & $\rho$.P$_0$.(1 + 2.$\rho$ + 3.$\rho^2$ + ...)\\
  & = & $\rho$.P$_0$/(1 - $\rho$)$^2$\\
  & = & $\rho$/(1 - $\rho$)  \hspace*{2in} since P$_0$ = (1 - $\rho$)\\
\end{tabular}

\noindent
At once

\begin{center}
\begin{tabular}{|| l c l  r ||}
\hline
         &&&\\
N        & = & $\rho$/(1 - $\rho$) &  \\
T        & = & E[s]/(1 - $\rho$)   &  \hspace*{1in}since T = N/$\lambda$ (Little's Law) \\
         &&&\\
T$_q$    & = & $\rho$.E[s]/(1 - $\rho$)  & since T = T$_q$ + E[s] \\
N$_q$    & = & $\rho^2$./(1 - $\rho$)  & since N$_q$ = $\lambda$.T$_q$
(Little's Law)\\
         &&&\\
\hline
\end{tabular}
\end{center}

\noindent
Finally, if we assume random arrival and service patterns,
then 
  
\begin{center}
\begin{tabular}{|| l c l r ||}
\hline
         &&&\\
$\pi_t$(r) & = & E[t].ln$\frac{100}{100-r}$ & \hspace*{0.5in}{\it r}th percentile value of T\\
           &   &  &\\
$\pi_t$(90) & $\simeq$ & 2.3 T        & = E[t].ln 10\\
$\pi_t$(95) & $\simeq$ & 3 T          & = E[t].ln 20\\
            &          &              & \\
$\pi_q$(r) & $\simeq$  & 
         $\frac{E[s]}{(1-\rho)}$.ln$\frac{100.\rho}{100-r}$ &\\[0.1in]
            & = &      $\frac{E[q]}{\rho}$.ln$\frac{100.\rho}{100-r}$ &\\
&&&\\
$\pi_q$(90) & $\simeq$ & E[t].ln(10.$\rho$) & \\
$\pi_q$(95) & $\simeq$ & E[t].ln(20.$\rho$) & \\
            &&&\\
\hline
\end{tabular}
\end{center}


\subsection*{Example 2}

A computer centre serves 20 randomly arriving customers 
per 16 hour day.
The customers have access to the computer via a single terminal.
The average time spent processing a query service is 30 minutes,
which is also randomly distributed.
The cpu is thus used 10 hours a working day, a utilisation of  0.625.
However the manager receives complaints about the long time customers have
to wait before logging on to the one terminal.
Are these complaints justified? How can this be with the terminal free
37.5\% of the time.

\noindent
{\bf Answer}.

We use the  M/M/1 model with $\rho$ = 0.625.

Then 

\begin{tabular}{ r c l c l r}
T & = & E[s]/(1-$\rho$) & = & 80 minutes & av. time in system   \\
T$_q$ & = & T - E[s]    & = & 50 minutes & average time in queue\\
N & = & $\rho$/(1-$\rho$) & = & 5/3      & (5/8)/(3/8)          \\
  & = & $\lambda$T        & = & 5/3      & (20/16).(4/3)        \\
N$_q$ & = & $\rho^2$/(1-$\rho$) & = & 1.041666.. & (5/8)$^2$/(3/8)\\      
      & = & $\lambda$T$_q$      & = & 1.041666.. & (20/16).(5/6)        \\
$\pi_q$(90) & = & T.ln(10$\rho$) & = & 146.61 mins & 90th percentile of times in queue \\
$\pi_t$(90) & $\simeq$ & 2.3T    & = & 184 mins & 90th percentile of times in system \\
\end{tabular}

\noindent
The complaints certainly seem justified. 
What can be done?

\noindent
\rule{160mm}{0.5mm} 


\newpage
\section{M/M/1/K --- Single server with a finite queue}

So far we have looked at the single server queue with
an infinite buffer. we now look at the M/M/1/K system 
in which a limit of K customers in the system is allowed.
When the system contains K customers, freshly arriving ones 
are turned away.

The arithmetic gets quite tricky, so let's look at some simple cases
first.


\subsection{M/M/1/2}

Here are the  M/M/1/2 system transitions:

\begin{figure}[hbt]\centering
\setlength{\unitlength}{0.16mm}
\begin{picture}(720,120)(0, 60)

\thicklines

\put(120,  120){\oval(60, 60)} 
\put(120,  120){\makebox(0, 0){\sf 0}}
\put(150,  140){\vector(1, 0){60}} 
\put(180,  150){\makebox(0, 0){$\lambda$}}


\put(240,  120){\oval(60, 60)} 
\put(240,  120){\makebox(0, 0){\sf 1}}
\put(270,  140){\vector(1, 0){60}} 
\put(300,  150){\makebox(0, 0){$\lambda$}}
\put(210,  100){\vector(-1, 0){60}} 
\put(180,   80){\makebox(0, 0){$\mu$}}

\put(360,  120){\oval(60, 60)} 
\put(360,  120){\makebox(0, 0){\sf 2}}
\put(330,  100){\vector(-1, 0){60}} 
\put(300,   80){\makebox(0, 0){$\mu$}}

\end{picture}
\label{M/M/1/2}
\end{figure}

\noindent
Let's derive some properties:

\begin{enumerate}

\item show that P$_0$, the probability of no customers in the system =
1/(1 + $\rho$ + $\rho^2$), and derive formulae for P$_1$ and P$_2$. 
P$_3$, P$_4$, ... etc are of course zero.

This system has only 3 states: S$_0$, S$_1$, and S$_2$.

P$_0$ + P$_1$ + P$_2$ = 1

The flow balance equations are 

\begin{tabular}{l c l ccc l c l}
$\lambda$.P$_1$ & = & $\mu$.P$_0$ 
   && or && P$_1$ & = & $\rho$.P$_0$\\
$\lambda$.P$_2$ & = & $\mu$.P$_1$ 
   && or && P$_2$ & = & $\rho^2$.P$_0$\\
\end{tabular}

Thus 

\begin{tabular}{l c l ccc l c l}
1  & = & P$_0$  & + & P$_1$ & + & P$_2$ \\
    & = & P$_0$  & + & $\rho$.P$_0$ & + & $\rho^2$.P$_0$ \\
    & = & \multicolumn{7}{l}{P$_0$(1 + $\rho$ + $\rho^2$)}\\[0.1in]
P$_0$ & = & \multicolumn{7}{l}{1/(1 + $\rho$ + $\rho^2$)}\\
P$_1$ & = & \multicolumn{7}{l}{$\rho$/(1 + $\rho$ + $\rho^2$)}\\
P$_2$ & = & \multicolumn{7}{l}{$\rho^2$/(1 + $\rho$ + $\rho^2$)}\\
\end{tabular}

Examples: if $\rho$ = 0.1, 0.5, and 0.9 then
(P$_0$, P$_1$, P$_2$) triples are (1/1.11, 0.1/1.11, 0.01/1.11),
(1/1.75,0.5/1.75,0.25/1.75),
and (1/2.71,0.9/2.71,0.81/2.71) respectively.
The chances of the system being full and an arrival rejected 
are 0.01/1.11, 0.25/1.75, and 0.81/2.71 respectively
(approx 1\%, 10\%, and 30\%).


\item find formulae for N and N$_q$.

\begin{tabular}{l c l ccc l c l}
N  & = & 0.P$_0$  & + & 1.P$_1$ & + & 2.P$_2$ \\
    & = & \multicolumn{7}{l}{$\rho$.P$_0$ + 2.$\rho^2$.P$_0$} \\
    & = & \multicolumn{7}{l}{$\rho$.(1+2.$\rho$).P$_0$}  \\
    & = & \multicolumn{7}{l}{$\rho$.(1+2.$\rho$)/(1 + $\rho$ + $\rho^2$)}\\[0.1in]
\end{tabular}

N$_q$ = 0.P$_0$ + 0.P$_1$ + 1.P$_2$ = P$_2$ = $\rho^2$/(1 + $\rho$ + $\rho^2$)


\item derive a formula for T, the average customer through time 

By Little's Law, T = N/$\lambda$', where the true arrival rate is
$\lambda$' = $\lambda$.(1-P$_2$) 
(with probability P$_2$ an arrival is discarded) 
= $\lambda$.(1 + $\rho$)/(1 + $\rho$  + $\rho^2$).

This is nice as the complicated terms cancel and we are left with:

T = E[s].(1+2.$\rho$)/(1+$\rho$)


\end{enumerate}

\subsection*{Exercises}


\begin{enumerate}

\item Do the same for M/M/1/1.

\item .... and M/M/1/3

\end{enumerate}


\subsection{M/M/1/K}

The state transition diagram for the M/M/1/K model is
typified below, with K=4 for concreteness:


\begin{figure}[hbt]\centering
\setlength{\unitlength}{0.20mm}
\begin{picture}(720,240)(0, 0)

\thicklines

\put(   0,   0){\framebox(720, 240){}}

\put(120,  120){\oval(60, 60)} 
\put(120,  120){\makebox(0, 0){\sf 0}}
\put(150,  140){\vector(1, 0){60}} 
\put(180,  150){\makebox(0, 0){$\lambda$}}


\put(240,  120){\oval(60, 60)} 
\put(240,  120){\makebox(0, 0){\sf 1}}
\put(270,  140){\vector(1, 0){60}} 
\put(300,  150){\makebox(0, 0){$\lambda$}}
\put(210,  100){\vector(-1, 0){60}} 
\put(180,   90){\makebox(0, 0){$\mu$}}

\put(360,  120){\oval(60, 60)} 
\put(360,  120){\makebox(0, 0){\sf 2}}
\put(390,  140){\vector(1, 0){60}} 
\put(420,  150){\makebox(0, 0){$\lambda$}}
\put(330,  100){\vector(-1, 0){60}} 
\put(300,   90){\makebox(0, 0){$\mu$}}

\put(480,  120){\oval(60, 60)} 
\put(480,  120){\makebox(0, 0){\sf 3}}
\put(510,  140){\vector(1, 0){60}} 
\put(540,  150){\makebox(0, 0){$\lambda$}}
\put(450,  100){\vector(-1, 0){60}} 
\put(420,   90){\makebox(0, 0){$\mu$}}

\put(600,  120){\oval(60, 60)} 
\put(600,  120){\makebox(0, 0){\sf 4}}
\put(570,  100){\vector(-1, 0){60}} 
\put(540,   90){\makebox(0, 0){$\mu$}}

\end{picture}
\label{M/M/1/K}
\caption{M/M/1/K transitions, K=4}
\end{figure}

\noindent
In the steady state, the flow balance equations yield:


\begin{center}
P$_1$ = $\rho$.P$_0$ \hspace*{0.5in}
P$_2$ = $\rho$.P$_1$ \hspace*{0.5in}
P$_3$ = $\rho$.P$_2$ \hspace*{0.5in}
.....................
P$_{K}$ = $\rho$.P$_{K-1}$
\end{center}


\noindent
Thus we have P$_k$ = $\rho^k$P$_0$ (for  k = 0, 1, 2, ..., K):

\begin{tabular}{ r c l}
1 & = & P$_0$ + P$_1$ + P$_2$ + P$_3$ + .... + P$_K$ \\
  & = & P$_0$.(1 + $\rho$ + $\rho^2$ + $\rho^3$ + .... + $\rho^K$)\\
  & = & P$_0$.(1 - $\rho^{K+1}$)/(1 - $\rho$)  \hfill for $\rho$ $<$ 1  \\
  &   & \\
\end{tabular}

\noindent
>From which 

\begin{tabular}{||r c l l||}
\hline
&&&\\
\hspace*{1in}
P$_k$ & = & $\frac{(1 - \rho).\rho^k}{1 - \rho^{K+1}}$
                   & if $\lambda$ $\neq$ $\mu$ \\[0.2in]
      & = & 1/(K-1)& if $\lambda$ = $\mu$ 
            \hspace*{1in} \\
&&&\\
\hline
\end{tabular}

\subsection{Derivation of key system properties}

The computation of the expected number in the system 
is split into two cases, as below:

\begin{enumerate}

\item $\lambda$ $\neq$ $\mu$.

\begin{tabular}{ l c l }
N & = & 0.P$_0$ + 1.P$_1$ + 2.P$_2$ + ... + K.P$_K$ \\
  & = & 1.$\rho$.P$_0$ + 2.$\rho^2$.P$_0$ + ... + K.$\rho^K$.P$_0$ \\
  & = & P$_0$.(1.$\rho$ + 2.$\rho^2$ + ... + K.$\rho^K$)\\
  & = & P$_0$.$\rho$.(1 + 2.$\rho$ + 3.$\rho^2$ + ... + K.$\rho^{K-1}$)\\
  & = & .... long and tedious, trust me ... \\
  & = & $\frac{\rho}{1-\rho}$ - $\frac{(K+1)\rho^{K+1}}{1-\rho^{K+1}}$\\
\end{tabular}

  
\item $\lambda$ = $\mu$.

\begin{tabular}{ l c l }
N & = & 0.P$_0$ + 1.P$_1$ + 2.P$_2$ + ... + K.P$_K$ \\
  & = & $\frac{1}{K+1}$(1 + 2 + ... + K) \\
  & = & $\frac{1}{K+1}$$\frac{K.(K+1)}{2}$\\
  & = & K/2\\
\end{tabular}

\end{enumerate}


In either case, N$_q$ is easy to compute.
We can compute N, the number in the system.
The server is idle with probability P$_0$,
in which case N = 0 and N$_q$ = 0.
The server is busy with probability (1 - P$_0$),
in which case N$_q$ is N less one (the guy being served).
Thus

\begin{center}
N$_q$ = N - (0.P$_0$ + 1.(1 - P$_0$)) = N - (1 - P$_0$)
\end{center}

Not all the traffic arriving at th system is permitted to
enter since the buffer has finite length (K-1).
An arrival is rejected if it arrives when the system is full,
which happens with probability P$_K$
Thus the effective arrival rate into the system is 

\begin{center}
$\lambda_a$ = $\lambda$.(1 - P$_K$)
\end{center}

\noindent
and the true system utilisation is $\rho$ = $\lambda_a$/$\mu$ rather 
than $\rho$ = $\lambda$/$\mu$.

By Little's Law,

\begin{center}
T = N/$\lambda_a$  \hspace*{1in} T$_q$ = N$_q$/$\lambda_a$
\end{center}


\subsection*{Example 3}

Traffic to a message switching center arrives randomly at an
average rate of 4 messages/second.
The messages take 0.22 seconds to transmit, randomly distributed.
Calculate the principal characteristics of the system 
given (i) an infinite number of buffers and (ii) 25 buffers.

\noindent
{\bf Answer}.

\begin{enumerate}


\item {\bf M/M/1 model}

The server utilisation, $\rho$ = $\lambda$.E[s] = 4$\times$0.22 =
 0.88.

\begin{tabular}{ r c l c l r}
T     & = & E[s]/(1-$\rho$) & = & 1.83 seconds & av. time in system   \\
T$_q$ & = & T - E[s]        & = & 1.61 seconds & average time in queue\\
N     & = & $\rho$/(1-$\rho$) & = & 7.33 jobs       & 0.88/0.12        \\
      & = & $\lambda$T        & = & 7.32 jobs      & 4.(1.83)          \\
N$_q$ & = & $\rho^2$/(1-$\rho$) & = & 6.45 jobs     & 0.88$^2$/0.12    \\      
      & = & $\lambda$T$_q$      & = & 6.44 jobs      & 4*1.61          \\
$\pi_q$(90) & = & T.ln(10$\rho$) & = & 3.98 seconds & 90th percentile of times in queue \\
$\pi_t$(90) & $\simeq$ & 2.3T    & = & 4.21 seconds & 90th percentile of times in system \\
\end{tabular}

The probability of 10 or more messages waiting 
(at least 11 in the system) is $\rho^{11}$ = 0.245, since

\begin{center}
\begin{tabular}{c l}
\multicolumn{2}{l}{P$_{11}$ + P$_{12}$ + P$_{13}$ + ..}\\
  = & $\rho^{11}$P$_0$ + $\rho^{12}$P$_0$ + $\rho^{13}$P$_0$ + ... )\\
  = & $\rho^{11}$P$_0$.(1 + $\rho$ + $\rho^{2}$ + ...)\\
  = & $\rho^{11}$\\
\end{tabular}
\end{center}

   



\item {\bf M/M/1/K model}

With the M/M/1/K model, messages arriving when the system is full are
discarded.
We would like to keep this percentage low, say less than half of
one percent (probability $<$ 0.005).
How many buffers are required?

With K-1 buffers, the probability that all are filled is

\begin{center}
P$_K$ = $\frac{(1-\rho).\rho^K}{(1-\rho^{K+1})}$, where $\rho$ = 0.88
\end{center}

We proceed by trial and error: P$_{20}$ = 0.00998. Too big.
P$_{25}$ = 0.0051. Still too big.
P$_{26}$ = 0.004464 which is acceptable --- 25 buffers.

Hence ($\lambda$ $\neq$ $\mu$)

\begin{tabular}{ r c l c l r}
$\lambda_a$  & = & (1-P$_k$).$\lambda$ & = & 3.982144 & effective
arrival rate\\
$\rho$       & = & $\lambda_a$/$\mu$   & = & 0.876072 & effective
service rate\\
N     & = &                   & = & 6.45 messages   & mean system size    \\
N$_q$ & = & N - (1-P$_0$)     & = & 5.57 messages   & mean messages waiting  \\
T     & = & N/$\lambda_a$     & = & 1.62 seconds    & av. time in system   \\
T$_q$ & = & N$_q$/$\lambda_a$ & = & 1.40 seconds    & average time in queue\\
\end{tabular}

All these numbers are smaller than those predicted by the M/M/1 model, 
which is as it should be.

\end{enumerate}

\noindent
\rule{160mm}{0.5mm} 

\newpage

\section{M/M/c --- multiple servers, infinite queue}

We now turn our attention to several servers per queue.
The mathematics gets quite complicated in a hurry,
so I will only present the argument and formulae
for the case of two servers, c = 2

The state transition diagram for the M/M/2 is shown
below. Notice that the arrival rates remain the same throughout.
With 2 or more customers in the system, both servers will be in
action, and then the expected completion rate is effectively doubled.

\begin{figure}[hbt]\centering
\setlength{\unitlength}{0.20mm}
\begin{picture}(720,240)(0, 0)

\thicklines

\put(   0,   0){\framebox(720, 240){}}

\put(120,  120){\oval(60, 60)} 
\put(120,  120){\makebox(0, 0){\sf 0}}
\put(150,  140){\vector(1, 0){60}} 
\put(180,  150){\makebox(0, 0){$\lambda$}}


\put(240,  120){\oval(60, 60)} 
\put(240,  120){\makebox(0, 0){\sf 1}}
\put(270,  140){\vector(1, 0){60}} 
\put(300,  150){\makebox(0, 0){$\lambda$}}
\put(210,  100){\vector(-1, 0){60}} 
\put(180,   88){\makebox(0, 0){$\mu$}}

\put(360,  120){\oval(60, 60)} 
\put(360,  120){\makebox(0, 0){\sf 2}}
\put(390,  140){\vector(1, 0){60}} 
\put(420,  150){\makebox(0, 0){$\lambda$}}
\put(330,  100){\vector(-1, 0){60}} 
\put(300,   88){\makebox(0, 0){2.$\mu$}}

\put(480,  120){\oval(60, 60)} 
\put(480,  120){\makebox(0, 0){\sf 3}}
\put(510,  140){\vector(1, 0){60}} 
\put(540,  150){\makebox(0, 0){$\lambda$}}
\put(450,  100){\vector(-1, 0){60}} 
\put(420,   88){\makebox(0, 0){2.$\mu$}}

\put(600,  120){\makebox(0, 0){\sf ...................}}

\put(570,  100){\vector(-1, 0){60}} 
\put(540,   88){\makebox(0, 0){2.$\mu$}}

\end{picture}
\label{M/M/2}
\caption{M/M/2 transitions}
\end{figure}

\noindent
In the steady state, the flow balance equations are:


\begin{tabular}{r c l r}
$\lambda$.P$_0$ & = & $\mu$.P$_1$ & flow balance between S$_0$ and
S$_1$ \\
&&&\\
$\lambda$.P$_1$ & = & $2.\mu$.P$_2$ & flow balance between S$_1$ and
S$_2$ \\
&&&\\
$\lambda$.P$_2$ & = & $2.\mu$.P$_3$ & flow balance between S$_2$ and
S$_3$ \\
&&&\\
\multicolumn{3}{c}{.................}\\
&&&\\
$\lambda$.P$_K$ & = & $2.\mu$.P$_{K+1}$ & flow balance between S$_K$ and
S$_{K+1}$, K $\geq$ 2 \\
&&&\\
\multicolumn{3}{c}{.................}\\
&&&\\
\end{tabular}


\noindent
With {\em $\rho$ = $\lambda$.E[s]/2} (the effective service rate when flat
out), these simplify to


\begin{center}
P$_1$ = 2.$\rho$.P$_0$ \hspace*{0.25in}
P$_2$ = $\rho$.P$_1$ \hspace*{0.25in}
P$_3$ = $\rho$.P$_2$ \hspace*{0.25in}
.....................\hspace*{0.25in}
P$_{K}$ = $\rho$.P$_{K-1}$
\end{center}


\noindent
Thus 

\begin{tabular}{ r c l}
1 & = & P$_0$ + P$_1$ + P$_2$ + P$_3$ + .... + P$_k$ + ...\\
  & = & P$_0$.(1 + 2.$\rho$ + 2.$\rho^2$ + 2.$\rho^3$ + .... +
  2.$\rho^k$ + ...)\\
  & = & P$_0$.(-1 + 2 + 2.$\rho$ + 2.$\rho^2$ + 2.$\rho^3$ + .... +
  2.$\rho^k$ + ...)\\
  & = & P$_0$.(-1 + 2(1 + $\rho$ + $\rho^2$ + $\rho^3$ + .... +
  $\rho^k$ + ...)\\
  & = & P$_0$.(-1 + 2/(1 - $\rho$))\\
  & = & P$_0$.(1 + $\rho$)/(1 - $\rho$)\\
\end{tabular}

\noindent
>From which 

\begin{tabular}{||r c l ||}
\hline
&&\\
\hspace*{1in}
$\rho$ & = & $\lambda$.E[s]/2 \\
&&\\
P$_0$  & = & (1 - $\rho$)/(1 + $\rho$) \\
P$_k$  & = & 2.P$_0$.$\rho^k$, for k = 1, 2, 3, ...\\
       & = & 2.$\frac{1 - \rho}{1 + \rho}$.$\rho^k$, for k = 1, 2, 3, ...
            \hspace*{1in} \\
&&\\
\hline
\end{tabular}

\subsection{Derivation of key system properties}

The computation of N, the expected number of customers in the system,
is as below:

\begin{tabular}{ l c l }
$\rho$ & = & $\lambda$.E[s]/2 \\
&&\\
N & = & 0.P$_0$ + 1.P$_1$ + 2.P$_2$ + 3.P$_3$ + ... + k.P$_k$ \\
  & = & 1.2.$\rho$.P$_0$ + 2.2.$\rho^2$.P$_0$ + 3.2.$\rho^2$.P$_0$ + 
            ... + k.2.$\rho^k$.P$_0$ +  ... \\
  & = & P$_0$.2.$\rho$.(1 + 2.$\rho$ + 3.$\rho^2$ + ... + k.$\rho^{k-1}$ + ... )\\
  & = & P$_0$.2.$\rho$/(1- $\rho$)$^2$\\
  & = & 2.$\rho$.$\frac{1 - \rho}{1 + \rho}$/(1- $\rho$)$^2$\\
  & = & 2.$\rho$./(1- $\rho^2$)\\
\end{tabular}

Hence, by Little's Law, T = N/$\lambda$ = E[s]/(1 - $\rho^2$).


Further, T$_q$ = T - E[s] = $\rho^2$.E[s]/(1 - $\rho^2$).

By Little's Law again, N$_q$ = $\lambda$.T$_q$ = 2.$\rho^3$/(1 - $\rho^2$)

\subsection*{Exercises}

\begin{enumerate}

\item Show that the probability of both servers busy in an M/M/2
system is:  2.$\rho^2$/(1 + $\rho$)

\item Given the M/M/2/2 system transitions:

\begin{figure}[hbt]\centering
\setlength{\unitlength}{0.16mm}
\begin{picture}(240,120)(120, 60)

\thicklines

\put(120,  120){\oval(60, 60)} 
\put(120,  120){\makebox(0, 0){\sf 0}}
\put(150,  140){\vector(1, 0){60}} 
\put(180,  150){\makebox(0, 0){$\lambda$}}


\put(240,  120){\oval(60, 60)} 
\put(240,  120){\makebox(0, 0){\sf 1}}
\put(270,  140){\vector(1, 0){60}} 
\put(300,  150){\makebox(0, 0){$\lambda$}}
\put(210,  100){\vector(-1, 0){60}} 
\put(180,   80){\makebox(0, 0){$\mu$}}

\put(360,  120){\oval(60, 60)} 
\put(360,  120){\makebox(0, 0){\sf 2}}
\put(330,  100){\vector(-1, 0){60}} 
\put(300,   80){\makebox(0, 0){2*$\mu$}}

\end{picture}
\caption{\label{M/M/2/2}M/M/2/2 transitions}
\end{figure}

\begin{enumerate}

\item show that P$_0$, the probability of no customers in the system =
1/(1 + 2$\rho$ + 2$\rho^2$), where {\bf $\rho$ = $\lambda$/2$\mu$},
and derive formulae for P$_1$ and P$_2$.

\item show that N, the average number of customers in the
system, is given by N =  2$\rho$$\frac{1+2\rho}{1 + 2\rho + 2\rho^2}$


\item derive a formula for T, the average customer through time.
Why is this formula so simple?


\end{enumerate}


\end{enumerate}


\newpage
\subsection*{Example 4}

What is the better solution for (a) the customer
and (b) the servers for the three cases:

\begin{enumerate}
\item  a single server with an arrival rate of $\lambda$,
a service rate $\mu$ and an average service time E[s]

\item  two separate M/M/1 queueing systems 
which share the arrivals and servicing (no queue jumping).
Each has  an arrival rate of $\lambda$/2,
a service rate $\mu$/2 and an average service time 2$\times$E[s]

\item a single queue with two servers which share a single queue of
arrivals (M/M/2).
The arrival rate is $\lambda$.
Each server has a service rate $\mu$/2 and an average service 
time of 2$\times$E[s]

\end{enumerate}

Notice that total arrival  rate to each system is the same;
and that each system has the same servicing power.

\noindent
{\bf Answer}.

Tabulate and compare the appropriate formulae.

Write a program to carry out the computations for
a selection of utilisations.

Here is just one comparison for E[s] = 1.0/second and 
an arrival rate of 0.5 customers/second.


\begin{center}
\begin{tabular}{|| c | c c c ||}
\hline
              & M/M/1      & 2 by M/M/1   & M/M/2   \\
\hline
$\lambda$     &  0.50       & 0.25         & 0.50   \\
E[s]          &  1.00       & 2.00         & 2.00   \\ 
$\rho$        &  0.50       & 0.50         & 0.50   \\
              &             &              &        \\
T             &  2.00       & 4.00         & 2.66   \\
T$_q$         &  1.00       & 2.00         & 0.66   \\
              &             &              &        \\
N             &  1.00       & 2.00         & 1.33   \\
N$_q$         &  0.50       & 1.00         & 0.33    \\
&&&\\
$\pi_t$(90)   &  4.60       & 8.00         & ---    \\
$\pi_q$(90)   &  3.22       & 6.44         & ---    \\
\hline
\end{tabular}
\end{center}





\noindent
What do these results tell us in the cases:

(i) we can replace a server by one that is quicker (e.g. cpu)

(ii) we cannot replace a server by one that is quicker, but
may employ more servers (e.g. bank clerks).

\noindent
\hfill{\rule{160mm}{0.5mm}}

\subsection*{Example 5}

\begin{enumerate}
\item Consider four variants of the same theme. 

\begin{enumerate}

\item M/M/1 system (one queue, one server with
          mean service rate of 1 job/minute)

\item an M/M/1/2 system (at most 2 jobs in the system,
so the maximum queue length is 1,
one server with mean service rate of 1 job/minute)

\item M/M/2 system (one queue, two servers each with a
          mean service rate of 0.5 jobs/minute)

\item M/M/2/2 system (one queue of maximum capacity 1, 
    two servers each with a mean service rate of 0.5 jobs/minute)

\end{enumerate}

For arrival rates $\lambda$ of 0.1, 0.5, and 0.9 jobs per minute,
tabulate $\rho$, N, and T (the utilisation, 
expected number in the system, and
expected through times).
Any conclusions you can draw?


The formulae are collected below:

\begin{center}
\begin{tabular}{lc lc lc l}
{\bf MM1} & & $\rho$ = ${\lambda}$/${\mu}$ &&
          N  = $\rho$/(1-$\rho$)&&
          T =  E[s]/(1-$\rho$)\\[0.1in]

{\bf MM1/2} & & $\rho$ = ${\lambda}$/${\mu}$ &&
          N  = $\rho$(1+2$\rho$)/(1+$\rho$+$\rho^2$)&&
          T =  E[s](1+2$\rho$)/(1+$\rho$)\\[0.1in]

{\bf MM2} & & $\rho$ = ${\lambda}$/2${\mu}$ &&
          N  = 2$\rho$/(1-$\rho^2$)&&
          T =  E[s]/(1-$\rho^2$)\\[0.1in]

{\bf MM2/2} && $\rho$ = $\lambda$/2$\mu$&&
N     =  2$\rho$(1+2$\rho$)/(1+2$\rho$+2$\rho^2$) &&
T     =  E[s] \\[0.1in]
\end{tabular}
\end{center}

{\bf Don't forget that $\rho$ = $\lambda$/2$\mu$ (or $\lambda$E[s]/2)
for both MM2 and MM2/2}. 

Here are the results:


\small
\begin{verbatim}
MM1

lam    mu     rho    Es       N           T
0.1    1.0    0.1    1.0    0.1111      1.1111
0.5    1.0    0.5    1.0    1.0000      2.0000
0.9    1.0    0.9    1.0    9.0000     10.0000

MM1/2

lam    mu     rho    Es       N           T
0.1    1.0    0.1    1.0    0.1081      1.0909
0.5    1.0    0.5    1.0    0.5714      1.3333
0.9    1.0    0.9    1.0    0.9299      1.4737

MM2

lam    mu     rho    Es       N           T
0.1    0.5    0.1    2.0    0.2020      2.0202
0.5    0.5    0.5    2.0    1.3333      2.6667
0.9    0.5    0.9    2.0    9.4737     10.5263


MM22

lam    mu     rho    Es       N           T
0.1    0.5    0.1    2.0    0.1967       2.0
0.5    0.5    0.5    2.0    0.8000       2.0
0.9    0.5    0.9    2.0    1.1403       2.0
\end{verbatim}
\normalsize

Conclusions: 

MM1 is better than MM2 with half speed servers (aside:
there will always be a need for fast cpu's).
This carries over to MM1/2 better than MM2/2 with half speed servers.
MMc (infinite buffer) is a poor approximation to MMc/n (finite buffer).

Here is my code (in SML with l for $\lambda$, m for $\mu$, and r for $\rho$):

\small
\begin{verbatim}
fun mm1 l m  = let val r  = l/m
                   val Es = 1/m
                   val N  = r/(1-r)
                   val T  = Es/(1-r)
               in
                   (l, m, r, Es, N, T)
               end;

fun mm12 l m = let val r  = l/m
                   val Es = 1/m
                   val N  = r*(1+2*r)/(1+r+r*r)
                   val T  = Es*(1+2*r)/(1+r)
               in
                   (l, m, r, Es, N, T)
               end;

fun mm2 l m  = let val r  = l/(2*m)
                   val Es = 1/m
                   val N  = 2*r/(1-r*r)
                   val T  = Es/(1-r*r)
               in
                   (l, m, r, Es, N, T)
               end;

fun mm22 l m = let val r  = l/(2*m)
                   val Es = 1/m
                   val N = 2*r*(1+2*r)/(1+2*r+2*r*r)
                   val T = Es
               in
                   (l, m, r, Es, N, T)
               end;
\end{verbatim}
\normalsize

Typical calls: mm1 0.1 1.0; mm12 0.1 1.0; mm2 0.1 0.5; mm22 0.1 0.5;
\end{enumerate}

\noindent
\hfill{\rule{160mm}{0.5mm}}




\newpage
\section{Jackson's Theorem for queueing networks}

Suppose a queueing network consists of N nodes satisfying the
conditions:

\begin{itemize}

\item each node consists of an M/M/c queueing system
with each server at node k having an average service time of
$\mu_k$

\item customers arriving at node k from the outside,
arrive at an average rate of $\lambda_k$.
Customers may also arrive at node k from other nodes within
the system.

\item once served at node k, a customer instantly goes to another node
j with probability p$_{kj}$ or leaves the network with probability
1 - $\sum_{j=1}^{n}$ p$_{kj}$

\end{itemize}

For each node k = 1, 2, 3, ..., n, the average arrival rate is

\begin{center}
$\Lambda_k$ = $\lambda_k$ + $\sum_{j=1}^{n}$ P$_{jk}$$\Lambda_j$
\end{center}

Further, if we let p(k$_1$, k$_2$, k$_3$, ..., k$_n$) denote
the steady state probability that there are k$_j$ customers
at node j, and if $\Lambda_j$ $<$ c$_j$$\mu_j$, then


\begin{center}
p(k$_1$, k$_2$, k$_3$, ..., k$_n$) =
p$_1$(k$_1$).p$_2$(k$_2$)...p$_n$(k$_n$) 
\end{center}

\noindent
where p$_i$(k$_i$) is the steady state probability of 
their being k$_i$ customers at node i if it were being treated
as an independent M/M/c system  
with an average arrival rate of $\Lambda_i$,
an average service  rate of $\mu_i$ for each of the
c$_i$ servers.



Fear not! It isn't quite as fierce as it looks.

If all the nodes are M/M/1 then
Jackson allows us to describe each node as an independent single
server system and we have:

\begin{center}
\begin{tabular}{||l c l||}
\hline
&&\\
P(N$_1$, N$_2$, ...N$_m$) & = &
$\frac{\rho_1^{n_1}}{1-\rho_1}$ $\times$ $\frac{\rho_2^{n_2}}{1-\rho_2}$ $\times$ ... 
$\times$ $\frac{\rho_m^{n_m}}{1-\rho_m}$\\
&&\\
E[N$_1$+N$_2$+...+N$_m$] & = & 
$\frac{\rho_1}{1-\rho_1}$ + $\frac{\rho_2}{1-\rho_2}$ + ... 
+ $\frac{\rho_m}{1-\rho_m}$\\
&&\\
E[T] & = & $\frac{E[N_1+N_2+...+N_m]}{\lambda_1 + \lambda_2 + ... +
\lambda_m}$ \\
&&\\
\hline
\end{tabular}
\end{center}



\newpage
\subsection*{Example 6}

Jobs arrive at a cpu at a mean rate of 4 per second randomly
distributed, where they execute for an average of 0.22 seconds,
also randomly distributed.
The probability  of a job running to completion is p;
and thus 1-p of requiring a further burst on the cpu.

\begin{figure}[hbt]\centering
\setlength{\unitlength}{0.20mm}
\begin{picture}(720,300)(0, 0)

\thicklines

\put(0,      0){\framebox(720, 270){}}

\put(360,  120){\oval(60, 60)} 
\put(360,  120){\makebox(0, 0){\footnotesize \sf M/M/1}}

\put(210,  120){\vector(1, 0){120}}
\put(390,  120){\vector(1, 0){120}}
\put(450,  120){\circle*{8}}
\put(450,  120){\line(0, 1){90}}
\put(450,  210){\line(-1, 0){180}}
\put(270,  210){\vector(0, -1){90}}

\put(200,  120){\makebox(0, 0){$\lambda$}}

\put(300,  135){\makebox(0, 0){$\Lambda$}}


\put(480,  135){\makebox(0, 0){p}}
\put(360,  225){\makebox(0, 0){1-p}}

\end{picture}
\label{j1}
\caption{Jackson System: 1}
\end{figure}

\noindent
{\bf Answer}.

By Jackson's theorem, $\Lambda$ = $\lambda$ + (1-p).$\Lambda$;
that is, $\Lambda$ = $\lambda$/p.

Further $\rho$ = $\Lambda$.E[s] = $\lambda$.E[s]/p.


Take p = 0.99. Then

\begin{tabular}{l c l c l}
$\Lambda$ & = & $\lambda$/p & = & 4.0404 \\
$\rho$ & = & $\Lambda$.E[s] & = & 0.8889\\
&&\\
N & = & $\rho$/(1 - $\rho$) & = & 8 messages \\
&&\\
T & = & E[s]/(1 - $\rho$)   & = & 1.98 seconds \\
\end{tabular}

\noindent
\hfill{\rule{160mm}{0.5mm}}

\subsection*{Example 7}

\begin{figure}[hbt]\centering
\setlength{\unitlength}{0.20mm}
\begin{picture}(720,340)(0, -40)

\thicklines

\put(0,      0){\framebox(720, 270){}}

\put(240,  120){\oval(60, 60)} 
\put(240,  120){\makebox(0, 0){\sf cpu}}

\put(480,  120){\oval(60, 60)} 
\put(480,  120){\makebox(0, 0){\sf io}}

\put(  0,  120){\vector(1, 0){210}}
\put(270,  120){\vector(1, 0){180}}
\put(510,  120){\line(1, 0){120}}

\put(360,  120){\circle*{8}}
\put(360,  120){\vector(0, -1){120}}

\put(630,  120){\line(0, 1){90}}
\put(630,  210){\vector(-1, 0){270}}
\put(360,  210){\line(-1, 0){270}}
\put( 90,  210){\vector(0, -1){90}}

\put(-20,  120){\makebox(0, 0){$\lambda$}}
\put(150,  135){\makebox(0, 0){$\Lambda_1$}}
\put(570,  135){\makebox(0, 0){$\Lambda_2$}}

\put(380,   60){\makebox(0, 0){p}}
\put(420,   135){\makebox(0, 0){1-p}}

\put(360,   -20){\makebox(0, 0){also $\lambda$}}
\end{picture}
\label{j2}
\caption{Jackson System: 2}
\end{figure}

All distributions are assumed to be random.
On average, each job takes 4 seconds of cpu time, 
with an IO interrupt every 0.25 seconds and 
each IO takes 0.2 seconds
The computer is in use 10 hours a day and processes 8000 jobs.

Find general formulae for, and then specific values for:

\begin{enumerate}

\item the cpu utilisation and the IO utilisation;

\item  the average number of jobs in the cpu queue (including
the job under execution);

\item  the average number of jobs in the IO queue (including
the job under execution);

\item  the average number of jobs in the system; and

\item the average response time of the system.

\end{enumerate}



\noindent
{\bf Answer}.

By Jackson's theorem:

\begin{center}
\begin{tabular}{l c l c l }
$\Lambda_1$ & = & $\lambda$ + $\Lambda_2$\\
$\Lambda_2$ & = & (1-p).$\Lambda_1$ \\
&&\\
$\Lambda_1$ & = & $\lambda$ + (1-p).$\Lambda_1$   & = & $\lambda$/p\\
$\Lambda_2$ & = & $\frac{1-p}{p}$.$\lambda$ \\
&&\\
$\rho_1$ & = & $\Lambda_1$/$\mu_1$ \\
$\rho_2$ & = & $\Lambda_2$/$\mu_2$ \\
&&\\
$\lambda_T$ & = & p.$\Lambda_1$                   & = & $\lambda$      \\
&&\\
p(n1, n2) & = & p(n1).p(n2) 
      & = & (1 - $\rho_1$)$\rho_1^{n_1}$.(1 - $\rho_2$)$\rho_2^{n_2}$.\\
&&\\
N     & = & $\rho_1$/(1 - $\rho_1$) +  $\rho_2$/(1 - $\rho_2$) \\
T     & = & N/$\lambda$ \\
\end{tabular}
\end{center}

\noindent
$\lambda$ = 8000 jobs per 10$\times$3600 seconds = 2/9 jobs/second.

The average number of IO interrupts per job is 4/0.25 = 16, so p = 1/16.

Hence $\Lambda_1$ = $\lambda$/p = 2/9$\times$16 = 32/9 jobs/sec;
and $\Lambda_2$ = $\frac{15}{16}$$\Lambda_1$ = $\frac{10}{3}$ jobs/sec.


\begin{enumerate}

\item $\rho_1$ = $\Lambda_1$.E[s$_1$] =
$\frac{32}{9}$$\times$$\frac{1}{4}$ = $\frac{8}{9}$, and
$\rho_2$ = $\Lambda_2$.E[2$_1$] =
$\frac{10}{3}$$\times$$\frac{1}{5}$ = $\frac{2}{3}$.

Hmmh. I wonder if 4/15 would be better for the 
expected time on the CPU making $\rho_1$= $\frac{128}{135}$?


\item The average number of jobs in the CPU system  is

E[N$_1$] = $\rho_1$/(1 - $\rho_1$) = 8 jobs; and
the average queue+service time at the CPU is 

E[T$_1$] = E[s$_1$]/(1 - $\rho_1$)  = 2.25 seconds.
Each job makes 16 passes through the CPU and thus spends 36 seconds in
that subsystem.

\item The average number of jobs in the IO system  is

E[N$_2$] = $\rho_2$/(1 - $\rho_2$) = 2 jobs
and the average queue+service time for IO is 
E[t$_2$] = E[s$_2$]/(1 - $\rho_2$)  = 0.6 seconds.
Each job makes 15 passes through the IO subsystem and thus spends 9
seconds there.

\item N    =  N$_1$ + N$_2$ = 10 jobs

\item The average response time is N/$\lambda$ =
10$\times$$\frac{9}{2}$ = 45 seconds.
This is made up of 4 seconds of CPU time and 3 seconds of IO time
--- the remaining 38 seconds are spent queueing.

\end{enumerate}


\noindent
\hfill{\rule{160mm}{0.5mm}}


\subsection*{Example 8}

 Consider a (Jackson) network consisting of three 
interconnected M/M/1 systems.

All external jobs entering the system do so at node 1.
Their arrival rate is 10 jobs/minute.

\begin{itemize}
\item from node 1, jobs go to node 2 with probability 0.4 
and to node 3 with probability 0.4,
and leave the network with probability 0.2 

The mean service rate at node 1, $\mu_1$, is 100 jobs/minute

\item from node 2, jobs go to node 1 with probability 1.0.
The mean service rate at node 2, $\mu_2$, is 80 jobs/minute


\item from node 3, jobs go to node 1 with probability 1.0
The mean service rate at node 3, $\mu_2$, is 25 jobs/minute
\end{itemize}


\begin{enumerate}
\item Draw a picture of the network showing the three nodes.

\item Write down the flow rate equations for each node,
and show that 
\begin{center}
$\Lambda_1$ = 50 jobs/minute; $\Lambda_2$ =  $\Lambda_3$ = 20 jobs/minute
\end{center}

where $\Lambda_1$, $\Lambda_2$, and $\Lambda_3$
are the actual flow rates at the 4 nodes respectively

\item What are the utilisations of the three nodes?

\item Which node is the bottleneck in this system?
At what flow rate of external jobs would the system saturate?

\end{enumerate}

\noindent
{\bf Solution}

\begin{enumerate}
\item Draw a picture of the network showing the three nodes.

\item  $\Lambda_1$ = 10 + $\Lambda_2$ + $\Lambda_3$

$\Lambda_2$ = 0.4.$\Lambda_1$ 

$\Lambda_3$ = 0.4.$\Lambda_1$ 

Hence $\Lambda_1$ = 10 + 0.4.$\Lambda_1$ + 0.4.$\Lambda_1$

Or 0.2.$\Lambda_1$ = 10

Hence $\Lambda_1$ = 50 jobs/minute; 
$\Lambda_2$ =  $\Lambda_3$ = 20 jobs/minute

\item What are the utilisations of the three nodes?

$\rho_1$ = 50/100 = 0.50; $\rho_2$ = 20/80 = 0.25; $\rho_3$ = 20/25 = 0.80

\item The bottleneck device is Node$_3$ at 80\%.
The utilisations rise with the external flow-in until 
the bottleneck device has 100\% utilisation.
That's Node$_3$ when the external flow rates has risen from 10
jobs/min to 10*100/80 = 12.5 jobs/minute.

\end{enumerate}


\noindent
\hfill{\rule{160mm}{0.5mm}}

\subsection*{Example 9}

Consider an open Jackson network comprising
three M/M/1 systems, N$_1$, N$_2$, N$_3$, with mean service rates 

\begin{center}
$\mu_1$ = 4 \hspace*{0.5in} $\mu_2$ = 4 \hspace*{0.5in} $\mu_3$ = 6
\end{center}

and mean customer arrival rates of

\begin{center}
$\lambda_1$ = 0.75 \hspace*{0.5in} 
$\lambda_2$ = 0.50 \hspace*{0.5in} 
$\lambda_3$ = 0.75
\end{center}

both measured in customers per minute.

On completion of service at N$_1$,
50\% of the jobs queue again at N$_1$,
25\% of the jobs queue at N$_2$, and
25\% of the jobs queue at N$_3$.

On completion of service at N$_2$,
25\% of the jobs queue at N$_1$ and
75\% of the jobs queue at N$_3$.

On completion of service at N$_3$,
50\% of the jobs queue again at N$_3$
and 50\% of the jobs leave the system.

\begin{enumerate}

\item Show that the effective flow rates at the three nodes are:


\begin{center}
$\Lambda_1$ = 2 \hspace*{0.5in} 
$\Lambda_2$ = 1 \hspace*{0.5in} 
$\Lambda_3$ = 4
\end{center}

measured in customers per minute.


The flow equations are :

\begin{tabular}{l c r c r c rc rc}
F1 & = & 0.75 & + & 0.50F1 & +  & 0.25F2 &  & \\
F2 & = & 0.50 & + & 0.25F1 &    &         &  & \\
F3 & = & 0.75 &   & 0.25F1 & +  & 0.75F2 & + & 0.50F3 \\
\end{tabular}

Take the equations for F1 and F2 and multiply through by 4:

\begin{tabular}{l c r c r c rc rc}
4F1 & = & 3 & + & 2F1 & +  & F2 &  & \\
4F2 & = & 2 & + & 1F1 &    &         &  & \\
\end{tabular}

\noindent
>From which F1 = 2 and F2 = 1. 
Substituting into the last equation for F3, we get F3 = 4.

Check: total flow in is 0.75 + 0.5 + 0.75 = 2 jobs/min.
Flow in = flow out = 0.5F3 = 2. Yes.


\item Hence find individual node utilisations 

Utilisations are 

$rho_1$ = 0.5; $rho_2$ = 0.25; $rho_3$ = 0.66; 

\item ... the expected number of customers at each node, and 
the expected number of customers in the system as a whole,

N = $\rho$/(1-$\rho$):


$N_1$ = 1.00; $rho_2$ = 0.33; $rho_3$ = 2.00; 

N in system = 3.33


\item ... and average time spent by a customer in the system.  


By Little's Law, T = N/($\lambda_1$ + $\lambda_2$ + $\lambda_3$) =
3.33/2 = 1.67 minutes.

\end{enumerate}



\noindent
\hfill{\rule{160mm}{0.5mm}}

\subsection*{Exercises}

\begin{enumerate}

\item Consider the 4-node queueing described below.
Jobs arrive externally at nodes 1 and 2 with equal rates
$\lambda$/2. The routeing is as follows:

\begin{itemize}
\item from node 1, jobs go to node 3 with probability 0.3 
and to node 4 with probability 0.7

\item from node 2, jobs go to node 3 with probability 0.7 
and to node 4 with probability 0.3

\item from node 3, jobs go back to node 1 with probability 0.8 
and leave the network with probability 0.2 

\item from node 4, jobs go back to node 2 with probability 0.8 
and leave the network with probability 0.2 

\end{itemize}

\begin{enumerate}

\item Draw a picture of the network showing the four nodes
and the flows amongst them

\item Write down the flow rate equations for each node

\item Show that

\begin{center}
$\Lambda_1$ + $\Lambda_2$ =  $\Lambda_3$ + $\Lambda_4$ = 5$\times$$\lambda$
\end{center}

where $\Lambda_1$, $\Lambda_2$, $\Lambda_3$, and $\Lambda_4$ 
are the actual flow rates at the 4 nodes respectively


\item ... and hence that $\Lambda_1$ = $\Lambda_2$ =  $\Lambda_3$ = $\Lambda_4$


\item If the average service times are all equal, give 
and justify an upper bound for the $\Lambda_k$.

\end{enumerate}
\end{enumerate}


\newpage
\section{Central Server Multiprogramming:\\Buzen's algorithm}

\begin{figure}[hbt]\centering
\setlength{\unitlength}{0.48mm}
\small
\begin{picture}(180, 200)(0, 0)

\thicklines

\put( -20,  0){\framebox(220, 200){}}

%wiring into/cpu%

\put( 40, 60){\framebox(20,80){\shortstack{T\\E\\R\\M\\I\\N\\A\\L\\S}}}
\put(100, 80){\framebox(40,40){SYSTEM}}

\put(  0, 100){\vector(1, 0){40}}
\put( 60, 100){\vector(1, 0){40}}
\put(140, 100){\line(1, 0){40}}
\put(180, 100){\line(0, 1){80}}
\put(180, 180){\line(-1, 0){80}}
\put( 80, 170){\makebox(20,20){$\lambda_T$}}
\put( 80, 180){\line(-1, 0){80}}
\put(  0, 180){\line(0, -1){80}}

\put( 50,   40){\makebox(0, 0){{\it THINK}}}
\put(120,   40){\makebox(0, 0){{\it RESPOND}}}

\put(40,  30){\vector(1, 0){20}}
\put(60,  30){\vector(-1, 0){20}}

\put(100,  30){\vector(1, 0){40}}
\put(140,  30){\vector(-1, 0){40}}


\put( 50,   20){\makebox(0, 0){{\it E[t]}}}
\put(120,   20){\makebox(0, 0){{\it T}}}

\end{picture}
\label{b0}
\caption{Abstract view}
\end{figure}


Consider a terminal driven system with N busy terminals.
Each  user has the  general computation pattern of:

\begin{center}
think; calculate; {\bf repeat}
\end{center}


\noindent
We assume that, {\bf on average, K of these users are actively
computing,  and N--K are thinking}.
Only the K are contributing to keeping the central system busy.
We can think of each user cyclically thinking and then submitting a
fresh task. The rate at which a fresh task is submitted to the 
central server system is $\lambda_T$.

Buzen's algorithm can be used to  answer {\it What if} questions.
What if we put in more terminals?
What if the expected think time goes down?
If these values are given, and we can compute the 
flow rate of tasks into/out of the central server system, $\lambda_T$,
then we can compute the system response time
by application of Little's Law {\it T + E[t] = N/$\lambda_T$}.

When a task is submitted, we assume that it will
make several IO requests which are uniformly
distributed over its computation. Suppose that, on average,
it makes 5 requests to device IO$_2$ and 4 requests to IO$_3$,
that is 9 IO requests in total.
Denoting a burst on the cpu by {\it b}, the pattern of each task is
taken to be:

\begin{center}
b; io$_2$; b; io$_2$; b; io$_2$; b; io$_2$; b; io$_2$;
b; io$_3$; b; io$_3$; b; io$_3$; b; io$_3$; b
\end{center}

\noindent
I.e. it makes 10 trips to the cpu, 5 to device 2, and 4 to device 3.
After a burst on the cpu, the
routing probabilities are

\begin{center}
\begin{tabular}{l c l}
P$_1$ & 0.1 & end of task; think again \\
P$_2$ & 0.5 & io request for device 2\\
P$_3$ & 0.4 & io request for device 3
\end{tabular}
\end{center}

\noindent
Notice that a task always includes a tidying up burst at the end.

Here is a more detailed picture of Buzen's model:



\begin{figure}[hbt]\centering
\setlength{\unitlength}{0.15mm}
\begin{picture}(640,700)(0, 0)

\thicklines

\put(  0,    0){\framebox(600, 690){}}

\put( 80,  660){\makebox(0, 0){\sf N--K thinkers}}
\put(390,  668){\makebox(0, 0){\sf K circulating programs}} 
\put(390,  642){\makebox(0, 0){\sf 1/p$_1$ laps each}} 


%wiring into/cpu%

\put(200,  360){\oval(80, 60)} 
\put(200,  360){\makebox(0, 0){$\mu_1$}}
\put(240,  360){\line(1, 0){80}}

\put(180,  390){\line(0,  1){90}}
\put(180,  480){\line(-1, 0){140}}
\put( 40,  480){\line(0, -1){120}}
\put( 40,  360){\vector(1, 0){40}}
\put( 80,  280){\framebox(40,160){}}

\put( 80,  300){\line(1, 0){40}}
\put( 80,  320){\line(1, 0){40}}
\put( 80,  340){\line(1, 0){40}}
\put( 80,  360){\line(1, 0){40}}
\put( 80,  380){\line(1, 0){40}}
\put( 80,  400){\line(1, 0){40}}
\put( 80,  420){\line(1, 0){40}}


\put(120,  360){\vector(1, 0){40}}

\put(110,  500){\makebox(0, 0){p$1$}}

%wiring into IO units%

\put( 320, 120){\line(0, 1){360}}
\put( 320, 480){\vector(1, 0){80}}
\put( 320, 360){\vector(1, 0){80}}
\put( 320, 240){\vector(1, 0){80}}
\put( 320, 120){\vector(1, 0){80}}

\put( 360, 500){\makebox(0, 0){p$_2$}}
\put( 360, 380){\makebox(0, 0){p$_3$}}
\put( 360, 140){\makebox(0, 0){p$_M$}}

\put(200,   40){\makebox(0, 0){\sf CPU}}
\put(440,   40){\makebox(0, 0){\sf IO devices}}


\put(440,  480){\oval(80, 60)} 
\put(440,  480){\makebox(0, 0){\sf $\mu_2$}}

\put(440,  360){\oval(80, 60)} 
\put(440,  360){\makebox(0, 0){\sf $\mu_3$}}

\put(440,  240){\makebox(0,0){.......................}}

\put(440,  120){\oval(80, 60)} 
\put(440,  120){\makebox(0, 0){\sf $\mu_M$}}


\put( 480, 480){\line(1, 0){80}}
\put( 480, 360){\line(1, 0){80}}
\put( 480, 240){\line(1, 0){80}}
\put( 480, 120){\line(1, 0){80}}


\put( 560, 480){\circle*{6}}
\put( 560, 360){\circle*{6}}
\put( 560, 240){\circle*{6}}
\put( 560, 120){\circle*{6}}

\put( 560, 120){\line(0, 1){480}}
\put( 560, 600){\line(-1, 0){340}}
\put( 220, 600){\vector(0, -1){210}}


\end{picture}
\label{b1}
\caption{Central Server Model}
\end{figure}

\noindent
Notice that whilst the arrival rate to the CPU from the terminals is
$\lambda_T$, the actual arrival rate to the CPU ($\lambda_1$)
is greater because tasks circulate 1/p1 times. 
It is easy to see that the two flow rates are related
by $\lambda_T$ = p$_1$.$\lambda_1$.


A cpu burst may be completed by an io request
in which case the task compete for IO device 2..M with probability 
p$_2$, p$_3$, ..., p$_M$; or the task is completed
with probability p$_1$ and the user enters think mode again at 
her/his terminal.
We assume that movements between queues are instantaneous and
there is no blocking (the queues always ``have space'').

\noindent
{\bf Model data:} what will always be supplied are the 
service rates $\mu_1$, $\mu_2$, ..., $\mu_M$ and
the routing probabilities p$_1$, p$_2$, ... p$_M$.
Don't start the calculation without them!


Let x$_1$ = 1 and then define

\begin{center}
x$_k$ = $\mu_1.p_k/\mu_k$ for k = 2, 3, ..., K
\end{center}

\noindent
These ratios are not plucked out of a hat,
they have a physical significance. Since the flow out at
the CPU is $\lambda_1$,
then the flow into device 2 is p$_2$$\lambda_1$,
the flow into device 3 is p$_3$$\lambda_1$, 
.... the flow into device k is p$_k$$\lambda_1$.

Notice that the utilisation of the CPU is $\lambda_1$/$\mu_1$,
and the utilisation of device k is p$_k$.$\lambda_1$/$\mu_k$.
So that x$_k$ is the ratio of the utilisation of device
k to that of the CPU:

{\large
\begin{center}
\begin{tabular} { r c l}
$\frac{\rho_k}{\rho_1}$ & = & $\frac{\lambda_k/\mu_k}{\lambda_1/\mu_1}$ \\
                       & = & $\frac{p_k.\lambda_1/\mu_k}{\lambda_1/\mu_1}$ \\
                       & = & $\frac{p_k/\mu_k}{1/\mu_1}$ \\
                       & = & $\frac{p_k/\mu_k}{1/\mu_1}$ \\
                       & = & $\frac{p_k/\mu_1}{\mu_k}$ \\
\end{tabular}
\end{center}}



{\bf k} = (k$_1$, k$_2$, ..., k$_M$) is a handy measure of the 
state of the system. It denotes the number of tasks at each device
(being served or queueing for service).

Buzen has shown that 


\begin{center}
\begin{tabular}{l c l}
p(k$_1$, k$_2$, ..., k$_M$)  & = &
  $\frac{1}{G(K)}$$\Pi_{j=2}^{M}(x_j)^{k_j}$\\
where &&\\
G(K)   & = & $\sum_{k \epsilon (K, M)} \Pi_{j=2}^{M}(x_j)^{k_j}$\\
\end{tabular}
\end{center}

\noindent
i.e. G(K) is defined so that the probabilities sum to 1.
It looks quite fearsome, but Buzen has a neat algorithm for
computing the key quantities, namely


\begin{center}
\begin{tabular}{||l c l||}
\hline
$\rho_1$ & = & G(K-1)/G(K)   \\[0.2in]
$\rho_j$ & = & $\rho_1$$x_j$ \hspace*{1in} j = 2, 3, ...,M\\
\hline
\end{tabular}
\end{center}

\noindent
 $\lambda_T$, the throughput in tasks per unit time, is
given by:

\begin{center}
$\lambda_T$ = $\mu_1$.p$_1$.$\rho_1$ 
              (or p$_1$.$\lambda_1$ if we knew $\lambda_1$)
\end{center}

\noindent
Given the number of terminals, N, 
and E[t] the  estimated think time (time spent in ``the terminal queue''),
the average response time (time from submitting a request to getting
the ``answer'' back) is 

\begin{center}
T = N/$\lambda_T$ - E[t] = $\frac{N}{\mu_1.p_1.\rho_1}$ - E[t]
\end{center}

\noindent
Finally the expected number in the queue of each IO subsystem, E[k$_j$],
is 

\begin{center}
E[k$_j$] = $\sum_{k=1}^{K}$(x$_j$)$^k$.$\frac{G(K-k)}{G(K)}$
\end{center}

\noindent
Here is an effective algorithm:


\begin{enumerate}

\item Set x$_1$ = 1 and compute the constants  x$_2$ .. x$_M$ using
x$_j$ = $\mu_1$.$\frac{p_j}{\mu_j}$ for j = 2, 3, ..., M.


\item Draw a g-table with rows labelled 0, 1, ...., K and columns
labelled x$_1$, x$_2$, ..., x$_M$. Set
the leftmost column and the topmost row to all ones.
I.e.

\begin{center}
\begin{tabular}{l c l l}
g(k, 1) & = & 1 & for k = 0, 1, ..., K \\[0.2in]
g(0, m) & = & 1 & for m = 1, 2, ..., M.
\end{tabular}
\end{center}


\item Iterate row by row. For the kth row, compute and enter

\begin{center}
 g(k, m) = g(k,m-1) + x$_m$.g(k-1,m) for m = 2, 3, .., M
\end{center}

\begin{figure}[hbt]\centering
\setlength{\unitlength}{0.40mm}
\begin{picture}(180,210)(0, -30)

\thicklines

\put(  0,    0){\oval(60, 60)} 
\put(  0,    0){\makebox(0, 0){g(k, m-1)}}

\put(  30,   0){\vector(1, 0){60}}

\put( 120,    0){\oval(60, 60)} 
\put( 120,    0){\makebox(0, 0){\Huge +}}

\put( 120,  120){\oval(60, 60)} 
\put( 120,  120){\makebox(0, 0){g(k-1, m)}}
\put(  75,  120){\makebox(0, 0){x$_m$$\times$}}

\put( 120,   90){\vector(0, -1){60}}


\end{picture}
\label{b2}
\caption{Template for the computation of g(k, m)}
\end{figure}


\item G(n) = g(n, M). I.e. The G's we want occupy the rightmost column.

\end{enumerate}

Hence $\rho_1$ from which the other system measures may be derived.


\subsection*{Example 10}

Due to Buzen.
Consider the case when K = 4 and M = 3; 
$\mu_1$ = 1/28,  $\mu_2$ = 1/40,  $\mu_3$ = 1/280 tasks/millisecond;
p$_1$ = 0.1, p$_2$ = 0.7, and p$_3$ = 0.2 respectively.


\noindent
{\bf Answer}.


\noindent
{\bf Step 1:} We first compute the constants x$_j$, j = 1, 2, 3.

x$_1$ = 1.

x$_2$ = $\mu_1$.p$_2$/$\mu_2$ = $\frac{1}{28}$$\times$0.7/(1/40) = 1.

x$_3$ = $\mu_1$.p$_3$/$\mu_3$ = $\frac{1}{28}$$\times$0.2/(1/280) = 2.


\noindent
{\bf Step 2:} Now initialise the K $\times$ M table for the g's.

The framework is:

\begin{center}
\begin{tabular}{|| c |r  r r|| }
\hline
                  & \multicolumn{1}{c}{x$_1$}
                  & \multicolumn{1}{c}{x$_2$}
                  & \multicolumn{1}{c||}{x$_3$}\\

K$\setminus$x$_j$ & =      & =      & =    \\
                  & 1      & 1      & 2     \\
\hline
0                 &        &        &       \\
1                 &        &        &       \\
2                 &        &        &       \\
3                 &        &        &       \\
4                 &        &        &       \\
\hline
\end{tabular}
\end{center}

\noindent
and we fill in the leftmost column and the topmost row with ones.

\begin{center}
\begin{tabular}{|| c |r  r r|| }
\hline
                  & \multicolumn{1}{c}{x$_1$}
                  & \multicolumn{1}{c}{x$_2$}
                  & \multicolumn{1}{c||}{x$_3$}\\

K$\setminus$x$_j$ & =      & =      & =    \\
                  & 1      & 1      & 2     \\
\hline
0                 & 1      &  1     &  1     \\
1                 & 1      &        &       \\
2                 & 1      &        &       \\
3                 & 1      &        &       \\
4                 & 1      &        &       \\
\hline
\end{tabular}
\end{center}

\noindent
{\bf Step 3:}

We now apply the formula 
{\em g(k, m) = g(k,m-1) + x$_m$.g(k-1,m)}
working from left to right, row by row.

For row 2:

g(1, 2) = 1 + 1.1 = 2 and g(1,3) = 2 + 2.1 = 4

\begin{center}
\begin{tabular}{|| c |r  r r|| }
\hline
                  & \multicolumn{1}{c}{x$_1$}
                  & \multicolumn{1}{c}{x$_2$}
                  & \multicolumn{1}{c||}{x$_3$}\\

K$\setminus$x$_j$ & =      & =      & =    \\
                  & 1      & 1      & 2     \\
\hline
0                 & 1      &  1     &  1     \\
1                 & 1      &  {\bf 2}      & {\bf 4}  \\
2                 & 1      &        &       \\
3                 & 1      &        &       \\
4                 & 1      &        &       \\
\hline
\end{tabular}
\end{center}


\noindent
and so on giving


\begin{center}
\begin{tabular}{|| c |r  r r|| }
\hline
                  & \multicolumn{1}{c}{x$_1$}
                  & \multicolumn{1}{c}{x$_2$}
                  & \multicolumn{1}{c||}{x$_3$}\\

K$\setminus$x$_j$ & =      & =      & =    \\
                  & 1      & 1      & 2     \\
\hline
0                 & 1      &  1     &  1     \\
1                 & 1      &  2     &  4     \\
2                 & 1      &  3     &  11     \\
3                 & 1      &  4     &  26     \\
4                 & 1      &  5     &  57     \\
\hline
\end{tabular}
\end{center}


Thus we now have that:

\begin{center}
G(0) = 1, G(1) = 4, G(2) = 11, G(3) = 26, and G(4) = 57
\end{center}

Thus the cpu utilisation $\rho_1$ for 1, 2, 3, 4 levels of
multiprogramming are 1/4, 4/11, 11/26 and 26/57 respectively.

When K = 4

\begin{center}
\begin{tabular}{l c l c l}
$\rho_2$ & = &  x$_2$.$\rho_1$ & = & 26/57 \\
$\rho_3$ & = &  x$_3$.$\rho_1$ & = & 52/57 \\
&&&&\\
$\lambda_T$ &= & $\mu_1$.p$_1$.$\rho_1$ & = & 0.001629 tasks/millisec\\
            &&                          & or &  1.629 tasks/sec\\
&&&&\\
T           & = & N/$\lambda_T$ - E[t] & = & 2.277 secs\\
            &   & \multicolumn{3}{r}{average response time with N=20
and the think time E[t]=10 secs}\\
\end{tabular}
\end{center}

\noindent
\hfill{\rule{160mm}{0.5mm}}

\subsection*{Example 11}

Due to Reiser.
All times are in millions of CPU instructions
where the CPU executes 750,000 instructions a second.
The multiprogramming level K = 10.
There are two IO devices, D$_2$ and D$_3$.
The mean path length between IO operations to D$_2$ is m$_2$ = 0.005
million instructions and to D$_3$ is is m$_3$ = 0.01 million instuctions.
The mean path length of a task m$_1$ = 0.1 million instructions.
Thus after a burst on the cpu, the chances of going to
the various devices are 

\begin{center}
the cpu : the drum : the disk = 1 : 10 : 20
\end{center}

This gives task completion probabilities of 

\begin{center}
p$_1$ = 1/31, p$_2$ = 20/31 and p$_3$ = 10/31
\end{center}


The mean length of a CPU burst is 1/31 * 0.1 million instructions,
or 1/310 million instructions.


The average access time to D$_2$ = 0.03 secs and to D$_3$ = 0.1 secs.
Thus

\begin{center}
\begin{tabular}{l c l c r}
$\mu_1$ & = & 310$\times$0.75 & = & 232.5 tasks/sec\\
$\mu_2$ & = & 1/0.03   & = & 33.33 tasks/sec\\
$\mu_3$ & = & 1/0.10 & = & 10.00 tasks/sec\\
\end{tabular}
\end{center}



\noindent
{\bf Answer}.


\noindent
Here is the g-table:


\begin{center}
\begin{tabular}{|| c |r  r r|| }

\hline
                  & {x$_1$}
                  & {x$_2$}
                  & {x$_3$}\\

K$\setminus$x$_j$ & =      & =      & =    \\
                  & 1      & 4.5      & 7.5     \\
\hline
0                 & 1      &  1.00       &  1.00     \\
1                 & 1      &  5.50     &  13.00     \\
2                 & 1      &  25.75   &  123.25     \\
3                 & 1      &  116.88 &  1041.25     \\
4                 & 1      &  526.94   & 8336.31     \\
5                 & 1      & 2372.22   & 64894.56\\
6                 & 1      & 10675.98  &497385.20\\
7                 & 1      &  48042.93 & 3778431.95\\
8                 & 1      & 216194.18 & 28554433.84\\
9                 & 1      & 972874.83 & 215131128.60\\
10                & 1      & 4377937.72& 1617861403.00\\
\hline
\end{tabular}
\end{center}






\begin{center}
\begin{tabular}{l c l c l}
$\rho_1$    & =  & G(9)/G(10) & = & 0.133\\
$\rho_2$ & = &  x$_2$.$\rho_1$ & = & 0.598\\
$\rho_3$ & = &  x$_3$.$\rho_1$ & = & 0.997 (serious bottleneck) \\
&&&&\\
$\lambda_T$ &= & $\mu_1$.p$_1$.$\rho_1$ & = & 0.997 tasks/sec\\
&&&&\\
T           & = & N/$\lambda_T$ - E[t] & = & 10.054 secs\\
            &   & \multicolumn{3}{r}{average response time with N=20
and the think time E[t]=10 secs}\\
\end{tabular}
\end{center}


\noindent
\hfill{\rule{160mm}{0.5mm}}


\subsection*{Exercise}


\begin{enumerate}

\item Modify example 10 (due to Buzen)  by 
adding in a fourth device equivalent in power to device 3 
(the current bottleneck) 
and which shares its workload (i.e. P3  and P4 are both 0.1).

Modify this example again by 
adding in fourth/fifth/sixth devices
each equivalent to device 3 and which share its workload 
(i.e. P3, P4, P5, P6 are all 0.05).

Any comments on the results?


\end{enumerate}



\end{document}


